\section{Eulerian ideals and bounded exponents}
\label{app:eulerian-foundations}

We prove the coinvariant algebra statements used in
Section~\ref{sec:hook-action}.  The classical descent filtration and
the Gram-ideal lemma give a short proof of the complementary
annihilator identity inside the ordinary coinvariant algebra.

Put $P_n=K[x_1,\ldots,x_n]$, $p_s=\sum_i x_i^s$,
$I_n=(p_1,\ldots,p_n)$ and $C_n=P_n/I_n$.
For a label set $I$, ordered increasingly, write
\[
 x_I=\prod_{i\in I}x_i,\qquad
 \Delta_I=\det(x_i^{a-1})_{1\le a\le |I|,\ i\in I}.
\]
Thus the ideal $J_{n,d}$ of Section~\ref{sec:hook-action} is
generated by $g_I=x_I\Delta_I$ for $|I|=d$.
Also put $J_{n,0}=C_n$ and $J_{n,n}=0$.
For $d\ge1$, let
\[
 B_{n,d}=\operatorname{span}\{x^\alpha:0\le\alpha_i<d\},
 \qquad F_{n,d}=\operatorname{image}(B_{n,d}\longrightarrow C_n).
\]

\begin{proposition}[Eulerian complements and annihilators]
\label{prop:standalone-eulerian-foundations}
For $1\le d\le n$, the quotient projection induces a graded
$S_n$-equivariant isomorphism
\begin{equation}\label{eq:standalone-eulerian-complement}
 F_{n,d}\xrightarrow{\ \sim\ } C_n/J_{n,d}.
\end{equation}
Moreover, for $0\le d\le n$,
\begin{align}
 \Frob(J_{n,d};q)
 &=\sum_{\lambda\vdash n}s_\lambda
       \sum_{\substack{T\in\SYT(\lambda)\\\des(T)\ge d}}
                        q^{\maj(T)},
       \label{eq:standalone-eulerian-character}\\
 \operatorname{Ann}_{C_n}(J_{n,d})&=J_{n,n-d}.
       \label{eq:standalone-eulerian-annihilator}
\end{align}
On the coinvariant differential forms
\[
 \Omega_n=
 \frac{P_n\otimes\Lambda(\theta_1,\ldots,\theta_n)}
 {(p_1,\ldots,p_n,\dd p_1,\ldots,\dd p_n)},
 \qquad \delta_j=\sum_i x_i^j\partial_{\theta_i}\quad(j\ge1),
\]
the map $D_d=\delta_d\cdots\delta_1:\Omega_n^d\to C_n$
has image $J_{n,d}$ for $0\le d<n$; $D_0$ means the identity.
Its apolar adjoint is $\Theta_1\cdots\Theta_d$.
\end{proposition}

\begin{proof}
\emph{Determinants and straightening.}
For $1\le d\le n$, the two determinant presentations agree:
\begin{equation}\label{eq:standalone-determinant-presentations}
 I_n+(\Delta_I:|I|=d+1)
       =I_n+(x_I\Delta_I:|I|=d).
\end{equation}
Expansion along the constant row gives the inclusion from left
to right.  Conversely, append $(1,x_j,\ldots,x_j^d)^t$ to the
$d$ columns indexed by a fixed $I$ and sum the resulting
determinants over $j\in[n]$.  Each summand is zero or a signed
$(d+1)$-label Vandermonde.  The last column of the sum is
$(n,p_1,\ldots,p_d)^t$, so modulo $I_n$ the sum is
$\pm n x_I\Delta_I$.  This proves the reverse inclusion, and
also proves $g_{[n]}=0$ in $C_n$ when $d=n$.

We claim that every polynomial is congruent to an element of
$B_{n,d}$ modulo
$L_{n,d}=I_n+(\Delta_I:|I|=d+1)$.
Induct on $d$ for all $n$, and then descend on monomial support
for fixed $(n,d)$.  If $d=1$, the presentation above puts every
$x_i$ in $L_{n,1}$.  If $n\le d$, the Artin basis, with
exponents $0\le\alpha_i<i$, suffices.
Otherwise write a monomial of nonempty support $I$ as $x_I f$.
Apply the induction hypothesis for $(|I|,d-1)$ in $K[x_i:i\in I]$:
\[
 f=b+\sum_{s=1}^{|I|}p_s(I)h_s
       +\sum_{\substack{T\subseteq I\\|T|=d}}\Delta_T u_T,
 \qquad p_s(I)=\sum_{i\in I}x_i^s,
\]
where every exponent in $b$ is at most $d-2$ and all coefficients
use only variables in $I$.  After multiplication by $x_I$, the
first term is bounded as required and every determinant term
belongs to $L_{n,d}$ by
\eqref{eq:standalone-determinant-presentations}.  Finally,
\[
 x_Ip_s(I)h_s\equiv
       -\sum_{j\notin I}x_Ix_j^s h_s\pmod{I_n}.
\]
Each monomial on the right has larger support; for $I=[n]$ the
sum is zero.  The empty-support monomial is $1$.
This completes both inductions.  Taking homogeneous components
throughout proves that $F_{n,d}\to C_n/J_{n,d}$ is a graded
$S_n$-equivariant surjection.

\emph{The descent filtration.}
The classical descent straightening theorem identifies
\[
 F_{n,d}=\operatorname{span}
   \left\{b_w=\prod_{i\in\Des(w)}
                   (x_{w(1)}\cdots x_{w(i)}):\des(w)<d\right\}.
\]
These descent monomials belong to the Garsia--Stanton basis of
$C_n$, and the descent-representation refinement gives
\begin{equation}\label{eq:standalone-descent-filtration}
 \Frob(F_{n,d};q)=\sum_{\lambda\vdash n}s_\lambda
       \sum_{\substack{T\in\SYT(\lambda)\\\des(T)<d}}
                       q^{\maj(T)}.
\end{equation}
See \cite[Theorems~3.12 and~4.1]
{AdinBrentiRoichman2005} for the descent representations and
\cite[Section~4, Theorem~4.3 and the paragraph following it]
{GnedinGorinKerov2013} for the maximum-exponent realization.
These statements are defined
over $\mathbb Q$ and hence hold over $K$.
Put
\[
 r_d=\dim J_{n,d},\qquad
 R_d=\#\{w\in S_n:\des(w)\ge d\}.
\]
The surjection just proved gives $r_d\ge R_d$.

\emph{The Gram ideal and complementary dimensions.}
Newton's identities give $p_s=0$ in $C_n$ for every $s>0$.
The standard inner product is nondegenerate on
$W_n=\{(z_i)\in K^n:\sum_i z_i=0\}$, since $n\ne0$ in $K$.
Thus
\[
 v_s=(x_1^s,\ldots,x_n^s)\in C_n\otimes W_n,
 \qquad \langle v_s,v_t\rangle=p_{s+t}=0
 \quad(s,t\ge1).
\]
For $1\le d<n$, a generator $g_I$ of $J_{n,d}$ is a coordinate minor of
$v_1,\ldots,v_d$.  The product of such a minor and a generator
of $J_{n,n-d}$ is multilinear and alternating in two blocks of
vector arguments of total size $n>\dim W_n$.  Apply
Lemma~\ref{lem:gram} with the arguments independent, and then
substitute these prefix vectors.  Every Gram generator becomes
zero.  The polynomial ideal identity remains valid under this
substitution in $C_n$, even when vectors occur in both blocks.
Consequently
\begin{equation}\label{eq:standalone-prefix-orthogonality}
                 J_{n,d}J_{n,n-d}=0.
\end{equation}
The perfect Frobenius pairing of $C_n$ now gives
$r_d+r_{n-d}\le n!$.  On the other hand, complementing every
permutation value replaces its descent number by $n-1-\des(w)$,
so $R_d+R_{n-d}=n!$.  Hence
\[
 n!=R_d+R_{n-d}\le r_d+r_{n-d}\le n!,
\]
which forces $r_d=R_d$ and $r_{n-d}=R_{n-d}$.
The cases $d=0,n$ are immediate.

The graded surjection is therefore an isomorphism, proving
\eqref{eq:standalone-eulerian-complement}; subtracting
\eqref{eq:standalone-descent-filtration} from the ordinary
coinvariant character proves
\eqref{eq:standalone-eulerian-character}.
Equation~\eqref{eq:standalone-prefix-orthogonality} gives the
actual inclusion $J_{n,n-d}\subseteq\operatorname{Ann}(J_{n,d})$,
and the same dimension equality makes it an equality.

\emph{The prefix map.}
The contractions descend to $\Omega_n$, since
$\delta_j(\dd p_s)=s p_{s+j-1}=0$ in $C_n$.
Applying $D_d=\delta_d\cdots\delta_1$ to
$\theta_{i_1}\cdots\theta_{i_d}$ gives $g_I$.
The map is $C_n$-linear, so its image is exactly $J_{n,d}$.
Under the standard apolar pairing, $\delta_j^*=\Theta_j$;
reversing the composition order gives
$D_d^*=\Theta_1\cdots\Theta_d$.
\end{proof}

\subsection{A finite certificate for the Gram lemma}

Let $Z$ be a nondegenerate quadratic space of dimension $N$. On
$\Lambda^dZ\otimes\Lambda^eZ$, use the raising and contraction
operators $L,c$ with the ordinary tensor signs of Lemma~\ref{lem:gram}.
For $t=d+e-N>0$, the operator
\begin{equation}\label{eq:gram-finite-certificate}
 R_{d,e}=\sum_{j\ge0}
 \frac{(-1)^j c^{j+1}L^j}
      {(j+1)!\,t(t+1)\cdots(t+j)}
 \qquad\text{satisfies}\qquad LR_{d,e}=\operatorname{id}.
\end{equation}
Indeed,
$Lc^{j+1}L^j=c^{j+1}L^{j+1}+(j+1)(t+j)c^jL^j$,
so the sum telescopes; it is finite because exterior raising is
nilpotent. Interpreting its terms as alternating multilinear
polynomials gives an actual certificate in the ideal of
cross-block inner products.
